\documentclass[12pt]{article}
\usepackage[utf8]{inputenc}
\usepackage{amsmath}
\usepackage{geometry}

\geometry{a4paper, margin=1in}

\begin{document}

\section*{Estimasi Nilai Rata-Rata Fungsi}

Diketahui fungsi $f(x) = \frac{2}{x}$ pada interval $[a, b] = [1, 9]$. Kita akan mengestimasi nilai rata-rata fungsi menggunakan jumlah terbatas dengan $n = 4$ sub-interval dan titik tengah (\textit{midpoint}).

\subsection*{1. Menentukan Lebar Sub-interval ($\Delta x$)}
Lebar masing-masing sub-interval adalah:
\[ \Delta x = \frac{b - a}{n} = \frac{9 - 1}{4} = \frac{8}{4} = 2 \]

\subsection*{2. Menentukan Titik Tengah ($m_i$)}
Interval $[1, 9]$ dibagi menjadi empat sub-interval: $[1, 3], [3, 5], [5, 7],$ dan $[7, 9]$. Titik tengah masing-masing adalah:
\begin{itemize}
    \item $m_1 = \frac{1+3}{2} = 2$
    \item $m_2 = \frac{3+5}{2} = 4$
    \item $m_3 = \frac{5+7}{2} = 6$
    \item $m_4 = \frac{7+9}{2} = 8$
\end{itemize}

\subsection*{3. Estimasi Nilai Rata-Rata ($\bar{f}$)}
Rumus estimasi nilai rata-rata dengan titik tengah adalah:
\[ \bar{f} \approx \frac{1}{b-a} \sum_{i=1}^{n} f(m_i) \Delta x = \frac{1}{n} \sum_{i=1}^{4} f(m_i) \]

Substitusi nilai fungsi $f(x) = \frac{2}{x}$:
\begin{align*}
\bar{f} &\approx \frac{1}{4} \left[ f(2) + f(4) + f(6) + f(8) \right] \\
\bar{f} &\approx \frac{1}{4} \left[ \frac{2}{2} + \frac{2}{4} + \frac{2}{6} + \frac{2}{8} \right] \\
\bar{f} &\approx \frac{1}{4} \left[ 1 + 0.5 + 0.333\dots + 0.25 \right] \\
\bar{f} &\approx \frac{1}{4} \left[ \frac{12}{12} + \frac{6}{12} + \frac{4}{12} + \frac{3}{12} \right] \times \frac{2}{2} \text{ (menggunakan penyebut 24)} \\
\bar{f} &\approx \frac{1}{4} \left[ \frac{24 + 12 + 8 + 6}{24} \right] = \frac{1}{4} \left[ \frac{50}{24} \right] = \frac{50}{96} = \frac{25}{48}
\end{align*}

\subsection*{Kesimpulan}
Estimasi nilai rata-rata fungsi $f(x)$ pada interval tersebut adalah:
\[ \bar{f} \approx \frac{25}{48} \approx 0.5208 \]

\end{document}